\section{Generic non-Bloch magic inheritance}
\label{sec:r7-inheritance}

Let $\Delta_j$ be the five anchor margins.  Choose a hopping range $r$ and a
buffer $b=L-\ell$ so that the explicit locality budgets
\begin{align}
 B_{\rm spec}&=8Mg_*^{-2}C_{\rm tail}e^{-\gamma r},\\
 B_{\rm prop}=B_{\rm ret}=B_{\rm coh}&=4\varepsilon_{\rm dyn}(b,T,q),\\
 B_{\rm res}&=2\eta\varepsilon_{\rm res}(b,\eta,q)
 \label{eq:r7-budgets}
\end{align}
satisfy $B_j<\Delta_j/2$.  This is always possible for fixed
$(\ell,\eta,T)$: take $r,b$ large and choose $q$ with
$\Phi(q)<\eta$.  Put
\begin{equation}
 \widehat\Delta=\min_j(\Delta_j-B_j)>0
\end{equation}
and define the observable Lipschitz constants
\begin{equation}
 K_{\rm spec}=4M/g_*^2,\quad
 K_{\rm prop}=K_{\rm ret}=K_{\rm coh}=2T,
 \quad K_{\rm res}=1/\eta,\qquad K=\max_jK_j.
 \label{eq:r7-Kobs}
\end{equation}

\begin{theorem}[Generic non-Bloch magic inheritance]
\label{thm:r7-inheritance}
At the anchor of Theorem~\ref{thm:r7-anchor}, let
\begin{equation}
 \delta_\theta=\frac{\widehat\Delta}{2KC_L},
 \qquad
 \delta_p=\frac{\widehat\Delta}{2KC_p}.
 \label{eq:r7-deltas}
\end{equation}
Then every $(\theta,p)$ satisfying
$|\theta|<\delta_\theta$ and $\|p-p_A\|<\delta_p$ obeys all five strict
operational magic inequalities.  Moreover
\begin{equation}
 I_{\rm magic}\setminus\Theta_C,
 \qquad I_{\rm magic}=(-\delta_\theta,\delta_\theta)
\end{equation}
is dense and conull in $I_{\rm magic}$.  For every angle in that set the
infinite Hamiltonian is bounded and self-adjoint, the magic observables are
intrinsic observables of that operator, and no exact finite common periodic
cover (hence no exact finite Bloch realization) exists.
\end{theorem}

\begin{proof}
Use finite-range truncation for the two spectral moments and finite-ball
truncation for the resolvent and dynamics.  Equations
\eqref{eq:r7-local-observable-lipschitz} and \eqref{eq:r7-dyn-lip} give the
$K_j(C_L|\delta\theta|+C_p\|\delta p\|)$ perturbation.  The two truncation
sides give $B_j$.  Equations~\eqref{eq:r7-deltas} keep their sum below every
anchor margin.  Thus the same parameter points remain magic.  R5 proves that
$\Theta_C$ is countable, dense, and null, while its complement is dense and
conull; it also proves that $\theta\notin\Theta_C$ admits no finite common
cover.  Intersecting with the open magic interval proves the final claims.
\end{proof}

\begin{corollary}[Arbitrarily large finite observation scale]
\label{cor:r7-any-scale}
For every finite $\ell$, fixed $\eta>0$, and fixed $T<\infty$, one can choose
a finite buffer $b$ and obtain $\delta_\theta(\ell)>0$.  Explicitly,
\begin{equation}
 \delta_\theta(\ell)
 \ge \frac{\widehat\Delta}{K\mu RC_{\rm Schur}}
 e^{-(\ell+b)/R}.
 \label{eq:r7-exponential-window}
\end{equation}
For fixed proof budget the angular window therefore scales as
$e^{-\ell/R}$ up to the displayed prefactor.  Exact global periodicity is
absent, but every prescribed finite scale has a nonzero conull-dense magic
window.
\end{corollary}
